%!TEX root = TheoretischePhysikC.tex


\section{Vorlesung: Zerfließen des Wellenpakets}
Der Ausdruck ,,Wellenpaket`` wird für räumlich begrenzte Funktionen verwendet, wie etwa $\Psi\in L^2(\mathbb{R}^3)$, nicht für ebene Wellen. Vom Standpunkt der statistischen Interpretation gibt es also nur solche ,,Pakete``, und die Diskussion ebener Wellen ist bestenfalls eine etwas verkürzte Art, über die Fourier-Transformation zu sprechen. In der Frühzeit der Quantenmechanik galten gut lokalisierte Wellenpakete außerdem als eine Art Ersatz für Teilchen. Gemessen
daran war eine Eigenschaft der Wellenpakete etwas störend: sie ,,flossen auseinander``,  was für Teilchen ja wirklich ein ungebührliches Betragen ist. Das Zerfließen des Wellenpakets galt danach als typisch quantenmechanischer Effekt. Wenn man allerdings $\Psi$ konsequent statistisch deutet, ist nichts Paradoxes an diesem Vorgang: warum soll schließlich eine Wahrscheinlichkeitsverteilung
nicht mit der Zeit breiter werden? Wir werden hier sehen, dass das Zerfließen der Wellenpakete gerade ein klassischer Effekt ist: er wird asymptotisch genau durch die Impulsverteilung bestimmt, wie bei einem Schwarm klassischer Teilchen.

Wir betrachten also jetzt Teilchen in $\mathbb{R}^3$ mit Hamiltonoperator
\begin{equation}
	H = \frac{\widehat{p}_\text{ges.}^2}{2m} = \frac{\widehat{p}_{x}^2}{2m}+\frac{\widehat{p}_{y}^2}{2m}+\frac{\widehat{p}_{z}^2}{2m}.
\end{equation}
In diesem Fall entspricht wie üblich in der QM $H$ der Energie eines freien Teilchens. 

Wir interessieren uns für die Wahrscheinlichkeitsverteilung des Ortes zu verschiedenen Zeiten, besonders zu großer Zeit, das heißt, für die Ergebnisse von
Ortsmessungen, wenn zwischen Präparation und Messung eine große Zeit $t$ verstrichen ist. Wir haben zwei äquivalente Alternativen, um dieses Problem anzugehen: das Heisenbergbild und das Schrödingerbild. Im Schrödingerbild würden wir die Schrödingergleichung 
\begin{equation}
	i\hbar \frac{\partial}{\partial t}\Psi = \widehat{H}\Psi
\end{equation}
lösen, und dann das Ergebnis als 
\begin{equation}
	p_t(x) \equiv |\Psi(x,t)|^2
\end{equation}
ablesen. Wir werden die Rechnung im Heisenbergbild durchführen, d.h. wir halten $\Psi$ (und dabei
die Dichtematix $\rho$) fest und berechnen die Operatoren
\begin{equation}
	\widehat{x}(t) \equiv e^{\frac{it}{\hbar} \widehat{H}}\widehat{x} e^{-\frac{it}{\hbar} \widehat{H}}, \quad \widehat{y}(t) \equiv e^{\frac{it}{\hbar} \widehat{H}}\widehat{y} e^{-\frac{it}{\hbar} \widehat{H}} \quad \text{und} \quad \widehat{y}(t) \equiv e^{\frac{it}{\hbar} \widehat{H}}\widehat{y} e^{-\frac{it}{\hbar} \widehat{H}}.
\end{equation}

Als nächstes suchen wir passende ,,Koordinaten``, in denen wir unsere Rechnung duchführen können. Wegen des Auftretens der Exponentialfunktion ist es naheliegend, eine Darstellug des Hilbertraums zu wählen, in der $H$ ein Multiplikationsoperator ist, also die Impulsdarstellung:
\begin{equation}
	\Phi(\mathbf{p})\equiv \Phi(p_x,p_y,p_z) \equiv \frac{1}{\sqrt{2\pi\hbar}} \iiint_{-\infty}^\infty \Psi(x)e^{-\frac{i}{\hbar}(p_x x+ p_y y+ p_z z)}\,dxdydz.	
\end{equation} 
Dann ist
\begin{equation}\label{eq:timedepx}
	\begin{split}
		(\widehat{x}(t)\Phi)(\mathbf{p}) &= e^{\frac{it}{\hbar} \frac{|\mathbf{p}|^2}{2m}}\left(\widehat{x}(t)e^{-\frac{it}{\hbar}\widehat{H}}\Phi\right)(\mathbf{p}) \\
		&= e^{\frac{it}{\hbar} \frac{|\mathbf{p}|^2}{2m}}\left(i\hbar \frac{\partial}{\partial p_x}\right)\left(e^{-\frac{it}{\hbar}\frac{|\mathbf{p}|^2}{2m}}\Phi(\mathbf{p})\right) \\
		&= t \frac{p_x}{m}\Phi(\mathbf{p}) + i\hbar \frac{\partial}{\partial p_x} \Phi(\mathbf{p}), \\ 
	\end{split}
\end{equation}
also
\begin{equation}
	\widehat{x}(t) = t \frac{\widehat{p}_x}{m} + \widehat{x}(0),
\end{equation}
wobei $\frac{\widehat{p}_x}{m}$ den \emph{Geschwindigkeitsoperator} (in der $x$-Richtung) darstellt. Es gelten ähnliche Ausdrücke für $\widehat{y}(t)$ und $\widehat{z}(t)$. Den Ausdruck $\frac{p_\mu}{m} = \frac{\partial}{\partial p_\mu} \frac{|\mathbf{p}|^2}{2m}$, $\mu = x,y,z$, bezeichnet man auch als \emph{Gruppengeschwindigkeit}, weil er die Geschwindigkeit von \emph{Gruppen}, also Wellenpaketen beschreibt.

Um die Zeitabhängigkeit des Erwartungswertes von $\widehat{x}(t)$ zu bekommen, braucht man nun nur noch auf beiden Seiten von (\ref{eq:timedepx}) die Spur mit $\rho =  |\Psi\rangle\langle\Psi|$ zu nehmen. Wir erhalten dabei
\begin{equation}
	\begin{split}
		\langle\widehat{x}\rangle(t) &= \tr(\rho \widehat{x}(t)) \\
		&= \tr\left(\rho \left(t \frac{\widehat{p}_x}{m} + \widehat{x}(0)\right)\right)\\
		&= t \frac{\langle\widehat{p}_x \rangle(0)}{m} + \langle\widehat{x}\rangle(0).
	\end{split}
\end{equation}

Die Operatorgleichung enthält aber viel mehr Information als nur das: Die Diagonalisierung der Operatoren $\widehat{x}(t)$ enthält ja alle Information über die Wahrscheinlichkeitsverteilungen des Ortes zur Zeit $t$ in beliebigen Zuständen. Das Aussehen der Ortswahrscheinlichkeitsverteilung zu großen Zeiten ist also völlig von der gemeinsamen Verteilung der Operatoren $\frac{\widehat{p}_\mu}{m}$ bestimmt. Am Verhalten von $\widehat{x}$, $\widehat{y}$ und $\widehat{z}$  ist also überhaupt kein Unterschied zwischen klassischen und quantenmechanischen Teilchen zu sehen.

Die interessanteren Unterschiede zwischen klassischer Mechanik und Quantenmechanik (und die wichtigeren Anwendungen der Quantenmechanik) finden sich allerdings bei Teilchen in äußeren Potentialen, oder in Wechselwirkung miteinander.